\documentclass[aps,prl,twocolumn,showpacs,superscriptaddress]{revtex4-2}
\usepackage{amsmath,amssymb,booktabs,xcolor,hyperref,amsthm}
\newtheorem{theorem}{Theorem}
\newtheorem{corollary}{Corollary}

\begin{document}

\title{Appendix KB2: d-Shell Hopf Symmetry Points --- Full Derivation\\
$S_H(d)$ for All Period 4--6 Transition Metals,
the Hopf Symmetry Energy, and the Gold Colour Derivation}

\author{Michel Robert Cabri\'e}
\affiliation{Independent Researcher, Victoria, Australia}
\email{ZeroFreeParameters@gmail.com}
\date{20 March 2026}

\begin{abstract}
This appendix derives $S_H(d)$ for all Period~4 transition metals, 
quantifies the Hopf Symmetry Energy stabilising half-filled and 
fully-filled d-shells, maps the full cross-period anomaly table, 
and derives gold's colour from d-shell closure combined with 
relativistic 6s contraction. Supporting Letter~83.
\end{abstract}

\maketitle

\section{$S_H(d)$ for the d-Shell}

The d-shell has capacity 10. With $n_d$ electrons:
\begin{equation}
    S_H(d) = \frac{\min(n_d,\,10-n_d)}{10}
\end{equation}

Maximum $S_H(d)=0.50$ at $n_d=5$ (half-fill). 
Minimum $S_H(d)=0$ at $n_d=0$ or $n_d=10$ (empty or closed).

\section{Period 4 Full Table}

\begin{center}
\small
\begin{tabular}{lccll}
\toprule
Sym & $Z$ & $n_d$ & $S_H(d)$ & BCT status \\
\midrule
Sc & 21 & 1  & 0.10 & Normal \\
Ti & 22 & 2  & 0.20 & Normal \\
V  & 23 & 3  & 0.30 & Normal \\
\textbf{Cr} & \textbf{24} & \textbf{5}  & \textbf{0.50} & \textbf{Half-fill attractor} \\
Mn & 25 & 5  & 0.50 & At sym.~pt., no promotion \\
Fe & 26 & 6  & 0.40 & Normal \\
Co & 27 & 7  & 0.30 & Normal \\
Ni & 28 & 8  & 0.20 & Normal \\
\textbf{Cu} & \textbf{29} & \textbf{10} & \textbf{0.00} & \textbf{Closed multiplet} \\
Zn & 30 & 10 & 0.00 & Closed, retains $4s^2$ \\
\bottomrule
\end{tabular}
\end{center}

\section{Hopf Symmetry Energy}

\begin{equation}
    E_{HS} = -\hbar\omega_{\rm OHC}\cdot\alpha_0^2\cdot S_H(d)\cdot 10
\end{equation}

Energy gain from $S_H(d)=0.40\to0.50$ (Cr anomaly):
\begin{equation}
    \Delta E_{HS}(\text{Cr}) = \hbar\omega_{\rm OHC}
    \cdot\alpha_0^2\cdot(0.50-0.40)\cdot10
    = \hbar\omega_{\rm OHC}\cdot5.5\times10^{-6}
\end{equation}

When $\Delta E_{HS} > E_{\rm promotion}(4s\to3d)$, the anomalous 
configuration is preferred. This crossover occurs at Cr and Cu 
in Period~4, Mo and Ag in Period~5, Au in Period~6.

\section{Cross-Period Table}

\begin{center}
\begin{tabular}{llll}
\toprule
Period & Half-fill & Full-fill & Confirmed \\
\midrule
4 & Cr: $3d^54s^1$ & Cu: $3d^{10}4s^1$ & \checkmark \\
5 & Mo: $4d^55s^1$ & Ag: $4d^{10}5s^1$ & \checkmark \\
6 & --- & Au: $5d^{10}6s^1$ & \checkmark \\
\bottomrule
\end{tabular}
\end{center}
Zero exceptions across three periods.

\section{The Gold Colour Derivation}

Au ($Z=79$): $5d^{10}6s^1$. The $5d$-shell achieves closure by 
the d-Shell Closure Theorem (Letter~83).

\textit{Consequence 1 --- Nobility:} Closed $5d$ multiplet 
$\to$ Noble Gas Theorem on d-shell $\to$ $5d$ electrons don't 
bond $\to$ gold is a noble metal.

\textit{Consequence 2 --- Colour:} The single $6s^1$ electron 
at $Z=79$ undergoes significant relativistic contraction, 
lowering its energy and shifting the $6s\to5d$ optical transition 
from ultraviolet into the visible blue ($\sim$450--480~nm). 
Gold absorbs blue light and reflects yellow-orange.

\begin{corollary}
Gold is gold-coloured because its $5d$-shell achieves BCT 
topological closure. The colour is a consequence of d-shell 
Hopf topology combined with relativistic $6s$ contraction at $Z=79$.
\end{corollary}

\begin{thebibliography}{99}
\bibitem{L83} M.~R.~Cabri\'e, \textit{BCT Letter 83}, Zenodo (2026).
\bibitem{L81} M.~R.~Cabri\'e, \textit{BCT Letter 81}, Zenodo (2026).
\bibitem{L82} M.~R.~Cabri\'e, \textit{BCT Letter 82}, Zenodo (2026).
\end{thebibliography}

\end{document}
